\chapter{Cơ sở lý thuyết và nguyên tắc bảo mật}
\label{chap:foundations}

Chương này trình bày các khái niệm được dùng để thiết kế và đánh giá nguyên
mẫu. Mục đích không phải thay thế một chuyên khảo an toàn thông tin, mà là làm
rõ vì sao từng thành phần của RBAC Guard tồn tại, dữ liệu nào nó xử lý và giới
hạn nào cần giữ khi diễn giải kết quả. Các nguyên tắc dưới đây được dùng trực
tiếp trong mô hình quyền, API Nova Bank, cấu trúc log và pipeline phát hiện.

\section{Kiểm soát truy cập theo vai trò}

Kiểm soát truy cập trả lời câu hỏi chủ thể nào được làm hành động nào trên tài
nguyên nào. Trong mô hình cấp quyền trực tiếp, quyền được gán vào từng tài khoản;
cách này đơn giản ở quy mô nhỏ nhưng nhanh chóng khó quản trị khi nhiều người có
cùng chức năng nghiệp vụ. RBAC giải quyết vấn đề bằng cách đưa vai trò vào giữa
user và permission: tổ chức gán quyền cho role, rồi gán role cho user. Mô hình
này được hình thức hóa bởi Ferraiolo và Kuhn, sau đó được Sandhu và cộng sự mở
rộng thành họ mô hình RBAC96 \citep{Ferraiolo1992RBAC,Sandhu1996RBAC}.

Trong RBAC lõi, tập user $U$, role $R$, permission $P$ và session $S$ được nối
bằng hai quan hệ gán: $UA \subseteq U \times R$ và $PA \subseteq P \times R$.
Permission trong nguyên mẫu được biểu diễn bằng cặp $(resource, action)$, ví
dụ \texttt{transactions:approve}. Một request chỉ được thực thi khi user active
có ít nhất một role dẫn đến permission yêu cầu. Cách biểu diễn này giúp kiểm
tra quyền nhất quán ở CLI, FastAPI và pipeline phân tích log.

Nguyên tắc \emph{đặc quyền tối thiểu} yêu cầu chỉ cấp những quyền cần cho nhiệm
vụ hiện tại. Vì vậy Teller được tạo giao dịch và xem các giao dịch thuộc phạm
vi của mình, còn Controller mới có quyền phê duyệt. Tách quyền như vậy giảm
xác suất một tài khoản đơn lẻ vừa khởi tạo vừa tự xác nhận giao dịch. RBAC Guard
chưa có role hierarchy, delegated administration hoặc review quyền định kỳ;
đây là các mở rộng hợp lý khi hệ thống chuyển từ demo sang vận hành thật.

\section{Xác thực, phân quyền và trách nhiệm giải trình}

Xác thực (authentication) xác định người đang gửi request; phân quyền
(authorization) xác định request đó có được phép không. Hai bước cần tách biệt:
một password hợp lệ chỉ tạo danh tính, không tự động cho phép mọi thao tác. API
Nova Bank thực hiện xác thực tại \texttt{/auth/login}, tạo token phiên, rồi
kiểm tra token và permission ở từng endpoint có bảo vệ. Mô hình này tránh sai
lầm phổ biến là để frontend tự quyết định một nút có được bấm hay không.

Sau khi request kết thúc, hệ thống cần accountability: khả năng trả lời ai đã
thực hiện gì, vào lúc nào, với kết quả nào. NIST SP 800-53 xem audit và
accountability là nhóm kiểm soát độc lập, nhấn mạnh việc tạo, bảo vệ và xem xét
bản ghi \citep{JTF2020SP80053}. Trong dự án, audit log lưu username, role tại
thời điểm thao tác, resource, action, outcome, thời điểm và transaction
reference. Nó không phải bằng chứng không thể sửa đổi như hệ thống pháp chứng,
nhưng đủ để minh họa quyết định allowed, denied, baseline bypass và conflict.

Nguyên tắc phân tách nhiệm vụ (separation of duties) bổ sung cho đặc quyền tối
thiểu. Nó không chỉ hỏi ``Teller có quyền hay không'' mà còn xác định các bước
nhạy cảm không nên do cùng một vai trò hoàn tất. Kịch bản Nova Bank giữ chế độ
baseline như một đối chứng có chủ đích: Teller có thể tạo và duyệt giao dịch.
Sau khi áp dụng RBAC, endpoint phê duyệt kiểm tra quyền tại backend và từ chối
Teller ngay cả khi người dùng biết URL của hồ sơ. Điều này thể hiện enforcement
và audit trail phải nằm tại điểm thực thi nghiệp vụ.

\section{Nhật ký bảo mật và vòng đời sự kiện}

Log bảo mật là dấu vết có cấu trúc của hoạt động xác thực, truy cập, thay đổi
quyền và lỗi hệ thống. Theo hướng dẫn quản lý log của NIST, log có giá trị cho
giám sát, điều tra sự cố và đáp ứng, nhưng chỉ hữu ích khi tổ chức có thể thu
thập, bảo quản và phân tích chúng nhất quán \citep{Kent2006SP80092}. Vấn đề
thực tế là các nguồn log thường không đồng nhất về cột, định dạng thời gian,
tên user và mức độ chi tiết.

OWASP khuyến nghị ứng dụng ghi nhận thành công/thất bại xác thực, lỗi phân
quyền và các thay đổi cấu hình có ý nghĩa bảo mật; đồng thời không đưa mật
khẩu, access token hoặc dữ liệu nhạy cảm không cần thiết vào log
\citep{OWASPLogging}. Hai yêu cầu này tạo một cân bằng: log phải đủ ngữ cảnh để
điều tra nhưng không trở thành kho dữ liệu làm tăng hậu quả khi bị lộ. Schema
Event của dự án vì thế ưu tiên định danh sự kiện, actor, nguồn, hành động, kết
quả và bằng chứng tối thiểu.

RBAC Guard dùng Event chuẩn hóa làm biên giữa dữ liệu đầu vào và detection
engine. Event giữ \texttt{event\_id}, timestamp, loại sự kiện, user, IP,
resource, action, status, request và details. Parser chấp nhận CSV/JSON, xác
nhận header và kiểu dữ liệu, sau đó cô lập RowError thay vì hủy toàn bộ lần
chạy vì một dòng lỗi. Thiết kế này phản ánh một nguyên tắc vận hành: dữ liệu
không hợp lệ phải có thể truy vết, trong khi dữ liệu hợp lệ vẫn có thể được
phân tích.

Vòng đời dữ liệu của nguyên mẫu gồm bốn bước. Thứ nhất, parser tạo Event có
thứ tự thời gian. Thứ hai, rule engine tạo finding với rule ID và evidence.
Thứ ba, scorer biến finding thành Alert có score và severity. Cuối cùng,
reporter ghi artifact để con người hoặc lệnh evaluate kiểm tra. Sự tách lớp này
giúp phân biệt lỗi dữ liệu, lỗi logic detection và lỗi báo cáo, đồng thời cho
phép tái chạy cùng input/configuration để kiểm chứng kết quả.

\section{SQL Injection và phát hiện chỉ báo trong log}

SQL Injection xảy ra khi dữ liệu không tin cậy được ghép vào câu lệnh truy vấn
theo cách làm thay đổi ý nghĩa truy vấn. Phòng vệ ưu tiên là prepared statement,
stored procedure phù hợp, kiểm tra allow-list và tài khoản database có quyền
tối thiểu, thay vì chỉ dựa vào filter chuỗi \citep{OWASPSQLi}. Vì vậy rule SQL
Injection trong RBAC Guard không được xem là lớp phòng vệ ứng dụng chính; nó là
lớp giám sát hậu kiểm để phát hiện request có chỉ báo đáng ngờ trong log.

Rule dùng một tập regex xác định trước và lưu pattern đã khớp làm evidence.
Ưu điểm của cách này là người đọc biết ngay vì sao Event bị gắn cờ, dễ đối
chiếu bằng chứng và tái lập cùng một kết quả. Đổi lại, nó có thể bỏ sót
encoding, payload phân
mảnh hoặc biến thể chưa nằm trong pattern; mở rộng pattern lại có thể bắt nhầm
văn bản hợp lệ. Kết quả SQL Injection của dataset tổng hợp vì thế được trình
bày như kiểm chứng implementation theo scenario, không phải tỷ lệ phát hiện
trên lưu lượng web thực.

Trong ngôn ngữ ATT\&CK, việc khai thác một ứng dụng public-facing có thể là
đường vào ban đầu của đối thủ \citep{MITRET1190}. Tuy nhiên một chuỗi xuất hiện
trong request log chỉ là dấu hiệu quan sát: nó không chứng minh truy vấn đã chạy
hoặc kẻ tấn công đã truy cập dữ liệu. Việc giữ phân biệt này giúp tránh diễn giải
một alert thành một incident đã được xác nhận.

\section{Password guessing và kiểm soát xác thực}

Password guessing hoặc brute force là chuỗi thử thông tin xác thực lặp lại để
tìm mật khẩu đúng. ATT\&CK mô tả brute force là kỹ thuật T1110, có thể nhắm vào
một tài khoản hoặc phân tán qua nhiều tài khoản, IP và khoảng thời gian
\citep{MITRET1110}. Kỹ thuật con T1110.001 tách password guessing khỏi password
spraying và credential stuffing \citep{MITRET1110001}. Trong dataset baseline,
nguyên mẫu chọn trường hợp đơn
giản có thể giải thích: nhóm login failure theo cặp $(user, IP)$ và báo động
khi số lần thất bại đạt ngưỡng trong một cửa sổ thời gian.

Ngưỡng không phải kết luận về ác ý. Một người dùng quên mật khẩu cũng có thể
tạo chuỗi failure, trong khi đối thủ có thể né ngưỡng bằng cách đổi IP, đổi user
hoặc kéo dài thời gian. NIST SP 800-63B cho phép dùng tín hiệu gian lận và rủi
ro như IP hoặc vị trí không mong đợi để quyết định kiểm soát xác thực bổ sung
\citep{NIST2025SP80063B}. Vì vậy context-risk của dự án thêm IP mới, hoạt động
ngoài giờ và denial lặp lại như bằng chứng triage, nhưng không tự khóa tài khoản
hay thay đổi quyền truy cập.

Một triển khai sản xuất có thể bổ sung rate limit, progressive delay, MFA,
credential-stuffing detection và cảnh báo phân tán. Những cơ chế này nằm ngoài
phạm vi PoC, nhưng giới hạn hiện tại cần được ghi rõ để người vận hành không
hiểu threshold năm lần/300 giây là chính sách áp dụng được cho mọi tổ chức.

\section{Chấm điểm rủi ro, ngữ cảnh và tương quan cảnh báo}

Kiểm tra quyền cho kết quả nhị phân, còn triage cần biết alert nào nên được xem
trước. Risk scoring kết hợp loại finding, confidence, số lần lặp và kết quả
RBAC để đưa về một thang 0--100. Score không phải xác suất tấn công được hiệu
chỉnh thống kê; nó là quy tắc ưu tiên minh bạch. Nhờ vậy người xem có thể lần
từ severity về điểm thành phần và evidence, thay vì phải tin vào một mô hình
hộp đen.

Tư duy zero trust nhấn mạnh việc đánh giá user, tài sản và tài nguyên thay vì
tin cậy chỉ vì vị trí mạng \citep{Rose2020SP800207}. Các tín hiệu ngoài giờ, IP
mới và tài nguyên hiếm trong dự án áp dụng ý tưởng này ở mức tối thiểu. Chúng
được học từ Event xuất hiện trước đó trong cùng lần chạy, nên phù hợp để minh
họa cách ngữ cảnh bổ sung cho rule nền, không phải hồ sơ hành vi dài hạn của
một khách hàng thực.

Một alert thường chỉ là mảnh nhỏ của chuỗi hành vi. Các nghiên cứu alert
correlation chỉ ra nhu cầu liên hệ alert để giảm gánh nặng phân tích
\citep{Valdes2001AlertCorrelation,Zhou2007AlertCorrelation}. RBAC Guard chọn
quy tắc gom nhóm đơn giản theo cùng user--IP và khoảng cách thời gian. Incident
kết quả giữ toàn bộ alert ID, Event ID, loại rủi ro, context signal và evidence;
điều này giúp người đọc truy ngược nguyên nhân. Quy tắc không mô hình hóa quan
hệ nhân quả như các hệ thống correlation chuyên sâu, đổi lại hành vi của nó dễ
xác minh và phù hợp với quy mô đồ án.
